%-------------------------
% MASTER RESUME -- everything, no page limit.
% Not for sending. Cut from this to build a targeted resume.
% Author : Kartikey Pathak
% Based on Jake's template
%------------------------

\documentclass[a4paper,11pt]{article}

\usepackage{latexsym}
\usepackage[empty]{fullpage}
\usepackage{titlesec}
\usepackage{marvosym}
\usepackage[usenames,dvipsnames]{color}
\usepackage{verbatim}
\usepackage{enumitem}
\usepackage[colorlinks=true, linkcolor=black, urlcolor=black, citecolor=black]{hyperref}
\usepackage{fancyhdr}
\usepackage[english]{babel}
\usepackage{tabularx}
\usepackage[normalem]{ulem}
\input{glyphtounicode}

\pagestyle{fancy}
\fancyhf{}
\fancyfoot{}
\renewcommand{\headrulewidth}{0pt}
\renewcommand{\footrulewidth}{0pt}

% Adjust margins
\addtolength{\oddsidemargin}{-0.6in}
\addtolength{\evensidemargin}{-0.6in}
\addtolength{\textwidth}{1.2in}
\addtolength{\topmargin}{-0.75in}
\addtolength{\textheight}{1.5in}

\urlstyle{same}

\raggedbottom
\raggedright
\setlength{\tabcolsep}{0in}

% Sections formatting
\titleformat{\section}{
  \vspace{-7pt}\scshape\raggedright\large
}{}{0em}{}[\color{black}\titlerule \vspace{-7pt}]

\pdfgentounicode=1

%-------------------------
% Custom commands
\newcommand{\resumeItem}[1]{
  \item\small{
    {#1 \vspace{-2pt}}
  }
}

\newcommand{\resumeSubheading}[4]{
  \vspace{-2pt}\item
    \begin{tabular*}{0.97\textwidth}[t]{l@{\extracolsep{\fill}}r}
      \textbf{#1} & #2 \\
      \textit{\small#3} & \textit{\small #4} \\
    \end{tabular*}\vspace{-7pt}
}

\newcommand{\resumeProjectHeading}[2]{
    \item
    \begin{tabular*}{0.97\textwidth}{l@{\extracolsep{\fill}}r}
      \small#1 & #2 \\
    \end{tabular*}\vspace{-7pt}
}

\newcommand{\resumeItemListStart}{\begin{itemize}}
\newcommand{\resumeItemListEnd}{\end{itemize}\vspace{-7pt}}
\newcommand{\resumeSubHeadingListStart}{\begin{itemize}[leftmargin=0.15in, label={}]}
\newcommand{\resumeSubHeadingListEnd}{\end{itemize}}

%-------------------------------------------

\begin{document}

%----------HEADING----------
\begin{center}
    \textbf{\Huge \scshape Kartikey Pathak} \\ \vspace{2pt}
    \small +91 7506048765 $|$
    \href{mailto:krpathak@proton.me}{krpathak@proton.me} $|$
    \href{https://linkedin.com/in/kartikey-pathak-a365a8281/}{\uline{LinkedIn}} $|$
    \href{https://github.com/NoobMaster-version}{\uline{GitHub}} $|$
    \href{https://noobmaster-version.github.io/spongefal/index.html}{\uline{Portfolio Website}}
    \\ \vspace{2pt}
    \textit{\footnotesize Master resume --- full record. Cut down per application.}
\end{center}

%-----------EDUCATION-----------
\section{Education}
  \resumeSubHeadingListStart
    \resumeSubheading
      {Veermata Jijabai Technological Institute (VJTI)}{Mumbai, India}
      {B.Tech in Electrical Engineering, Minor in Robotics (CGPA: 8.04)}{2023 -- 2027}
    \resumeSubheading
      {Matoshree Prabodhini Jr. College}{Maharashtra, India}
      {HSC, Maharashtra State Board --- MHT-CET 99.14 percentile}{2023}
  \resumeSubHeadingListEnd

%-----------EXPERIENCE-----------
\section{Experience}
  \resumeSubHeadingListStart

    \resumeSubheading
      {Eyecandy Robotics}{April 2026 -- August 2026}
      {Electronics and Hardware Engineer}{Bangalore, India}
      \resumeItemListStart
        \resumeItem{Designed production PCBs for a legged consumer robot: Raspberry Pi Zero 2W and Milk-V compute carriers, 1S/2S battery charging and power boards (BQ25887, MP2672, TPS562202 rails), a servo driver board, a microphone board and inter-board connector boards}
        \resumeItem{Owned schematic capture, routing, DFM review and fabrication output for each board}
        \resumeItem{Integrated MPU6050 IMU, a buffered Feetech servo bus, per-leg force-sensitive-resistor conditioning (LM393 comparators), PAM8403 audio and a ReSpeaker microphone array at board level}
        \resumeItem{Laid out test points and solder-jumper configuration options so each subsystem could be probed or bypassed independently during hardware bring-up}
        \resumeItem{Wrote fleet-provisioning tooling that clones one built Raspberry Pi runtime image onto many robots with unique machine IDs, hostnames and SSH host keys, so a fleet coexists on one network without identity or IP collisions}
      \resumeItemListEnd

    \resumeSubheading
      {NishCorp Technologies}{July 2025 -- November 2025}
      {Robotics Control Systems Engineer Intern}{Hybrid}
      \resumeItemListStart
        \resumeItem{Developed Gazebo simulation for a stair-climbing robotic wheelchair platform to analyze motion and steering behaviour, converting CAD to URDF and bringing the model up in Gazebo and RViz}
        \resumeItem{Applied differential drive kinematic models, integrated an IMU plugin, and wrote teleoperation so the platform could be driven and studied in simulation before it was built in steel}
        \resumeItem{Contributed to the half-scale hardware prototype through 3D printing, assembly, motor electronics, integration, testing and field evaluation; the prototype could sit and stand under its own power}
        \resumeItem{Simulation work formed part of the technical case that helped the startup raise Rs 13 lakh in grants}
      \resumeItemListEnd

    \resumeSubheading
      {\textbf{\href{https://noobmaster-version.github.io/spongefal/projects/kdail.html}{Krishna Defence and Allied Industries Limited}}}{April 2025 -- July 2025}
      {Research Intern - Defence Technology}{IIT Gandhinagar}
      \resumeItemListStart
        \resumeItem{Solved indoor localization for an autonomous forklift in a naval dockyard, where GPS is blocked by surrounding steel and wheel odometry drifts on every slip}
        \resumeItem{Surveyed Bluetooth, LoRa, RFID, radar and camera-based options before selecting ultra-wideband ranging}
        \resumeItem{Developed firmware for Qorvo DWM3001CDK Ultra-Wideband modules using Zephyr RTOS, handling low-level drivers and RF packets}
        \resumeItem{Implemented a custom double-sided two-way ranging protocol, tuning channel selection, preamble length and PAC size and re-deriving timing delays per configuration, extending usable range from the stock 45\,m to 315\,m}
        \resumeItem{Implemented SLAM and NAV2 on a Raspberry Pi 4B fusing UWB, IMU and LIDAR for autonomous forklift POC navigation in warehouse environments}
        \resumeItem{Scaled testing from a 2$\times$2\,m room to a 50$\times$60\,m yard and validated against a surveyor's total station rather than internal estimates: localization error of 30\,cm across a 10,000 sq meter area and 2\,cm in a 50 sq meter area}
      \resumeItemListEnd

  \resumeSubHeadingListEnd

%-----------PROJECTS-----------
\section{Projects}
  \resumeSubHeadingListStart

    \resumeProjectHeading
      {\textbf{PCB Axial-Flux BLDC Motor} $|$ \emph{KiCAD, FEMM, Python}}{Aug 2026 -- Present}
      \resumeItemListStart
        \resumeItem{Designing a coreless axial-flux motor whose stator is the PCB itself: 100\,mm two-layer board, 6 wedge coils over 8 poles, dual 3D-printed rotor carrying 16 N42 magnets across a 0.5\,mm air gap}
        \resumeItem{Generate the coil geometry programmatically rather than drawing it, so pole count, trace width and radii are parameters of the design}
        \resumeItem{Established that a top-down 2D cut cannot yield torque for an axial-flux machine, and built the magnetostatic model by unrolling the machine on a cylinder into five radial slices so that coil current runs out of plane}
        \resumeItem{Extracted torque constant, back-EMF, Kv and torque ripple from the solved air-gap flux, predicting 40.9\,mN$\cdot$m/A rms, 496\,rpm/V and zero cogging torque for the coreless design}
        \resumeItem{Identified before fabrication that the 3D-printed rotor gives away 3.3x the torque of the steel-backed design the board was sized for, and that the motor is thermally rather than magnetically limited; sized 6\,mm steel back plates as the fix by checking plate saturation}
      \resumeItemListEnd

    \resumeProjectHeading
      {\textbf{\href{https://github.com/NoobMaster-version/Field-Orient-Control-for-BLDC-Motor}{FOC Controller for Micro BLDC Motor}} $|$ \emph{STM32, DRV8311H, PCB Design}}{Jan 2026 -- Feb 2026}
      \resumeItemListStart
        \resumeItem{Developed a Field-Oriented Control (FOC) system for a 4300\,KV micro BLDC motor operating from a 3.7\,V supply}
        \resumeItem{Designed and routed a custom PCB integrating an STM32 microcontroller, DRV8311H motor driver, and AS5600 encoder}
        \resumeItem{Implemented centre-aligned three-phase PWM at 20\,kHz with a 40\,kHz control ISR, zero-sequence injection, AS5600 angle over I2C and DMA-fed phase current sensing on an STM32G431}
      \resumeItemListEnd

    \resumeProjectHeading
      {\textbf{Smart Switch Board} $|$ \emph{ESP32, ESP-IDF, FreeRTOS, MQTT, BLE, Flutter, PCB Design}}{2025 -- Present}
      \resumeItemListStart
        \resumeItem{Built a network-controlled mains wall switch across every layer alone: CAD enclosure, a perf board carrying relays and optocouplers to isolate 220\,V from a 3.3\,V ESP32, ESP-IDF firmware, a TLS MQTT broker and a Flutter application}
        \resumeItem{Held to a hard requirement that the physical switches keep working with network, broker and app all unavailable}
        \resumeItem{Implemented WiFi provisioning over BLE with no hardcoded credentials, switch state persisted in NVS across power cuts, and a 5-second heartbeat so switch, firmware, broker and app stay agreed on state}
        \resumeItem{Installed and running at home; next steps are a routed KiCAD PCB in place of the perf board and OTA firmware updates}
      \resumeItemListEnd

    \resumeProjectHeading
      {\textbf{\href{https://noobmaster-version.github.io/spongefal/projects/mario.html}{MARIO — 3DOF Robotic Manipulator}} $|$ \emph{ROS2, Gazebo Fortress, ESP32, micro-ROS}}{Apr 2025 -- Aug 2025}
      \resumeItemListStart
        \resumeItem{Developed embedded firmware in ESP-IDF and a micro-ROS communication pipeline for real-time hardware-software integration, so one command set drives both the simulation and the physical arm}
        \resumeItem{Migrated the simulation environment off end-of-life Gazebo Classic to Gazebo Fortress, porting model, plugins and launch files, and rewrote teleoperation for beginners}
        \resumeItem{Manufactured, tested and distributed 50+ robotics kits used by 200+ students for learning manipulator control}
      \resumeItemListEnd

    \resumeProjectHeading
      {\textbf{\href{https://noobmaster-version.github.io/spongefal/projects/kurma.html}{Kurma - Quadruped Robot}} $|$ \emph{ESP-IDF, MQTT, Kinematics}}{Jan 2025 -- Mar 2025}
      \resumeItemListStart
        \resumeItem{Designed and built a turtle-inspired quadruped robot, developing multiple gait patterns for stable multi-terrain locomotion}
        \resumeItem{Developed a teleoperation system utilizing MQTT for low-latency communication between a Raspberry Pi and ESP32s}
        \resumeItem{Used a statically stable gait with three feet always grounded, which kept the robot walking on rough ground that stops wheeled platforms}
      \resumeItemListEnd

    \resumeProjectHeading
      {\textbf{Open Manipulator X Simulation} $|$ \emph{ROS2, Gazebo, RViz, Computer Vision}}{June 2024 -- July 2024}
      \resumeItemListStart
        \resumeItem{Programmed forward and inverse kinematics for an Open Manipulator X robot arm to perform precise pick-and-place tasks}
        \resumeItem{Implemented object detection and tracking algorithms to locate target objects in the Gazebo simulation environment}
        \resumeItem{Chained speech recognition to object detection to inverse kinematics, so a spoken target became an arm trajectory; first project built with ROS2 nodes and topics rather than a single control loop}
      \resumeItemListEnd

    \resumeProjectHeading
      {\textbf{\href{https://github.com/NoobMaster-version/Coin_Detection}{Automated Coin Sorter}} $|$ \emph{Raspberry Pi, TensorFlow, OpenCV, Servos}}{24-Hour Hackathon}
      \resumeItemListStart
        \resumeItem{Classified Rs 1, 2, 5 and 10 coins from a camera feed using a CNN trained by transfer learning, reaching above 95\% accuracy on Rs 5 and Rs 10 coins}
        \resumeItem{Routed each classified coin into its bin using servo-actuated gates driven from the Raspberry Pi}
      \resumeItemListEnd

  \resumeSubHeadingListEnd

%-----------ADDITIONAL PROJECTS-----------
\section{Additional Projects}
  \resumeSubHeadingListStart

    \resumeProjectHeading
      {\textbf{\href{https://github.com/SRA-VJTI/Wall-E}{Self-Balancing Robot (Wall-E)}} $|$ \emph{ESP32, ESP-IDF, MPU6050, PID}}{2024}
      \resumeItemListStart
        \resumeItem{Balanced a two-wheeled robot by fusing gyroscope and accelerometer with a complementary filter and running a PID loop at a few hundred hertz; the platform SRA now teaches sensor fusion on}
      \resumeItemListEnd

    \resumeProjectHeading
      {\textbf{Maze Solving Robot} $|$ \emph{ESP32, ESP-IDF, IR Sensor Array}}{2024}
      \resumeItemListStart
        \resumeItem{Mapped a maze with the left-hand rule at a 100\,Hz control loop, recorded junctions, then replayed the shortest path; placed 4th in the SRA maze competition}
      \resumeItemListEnd

    \resumeProjectHeading
      {\textbf{\href{https://github.com/NoobMaster-version/gesture-controlled-car}{Gesture Controlled Car}} $|$ \emph{ESP32, MPU6050, ESP-NOW}}{Dec 2023 -- Jan 2024}
      \resumeItemListStart
        \resumeItem{Drove a car by hand tilt, reading a glove-mounted IMU on one ESP32 and transmitting over ESP-NOW to a second ESP32 running the motor drivers; built during a first-semester winter internship}
      \resumeItemListEnd

    \resumeProjectHeading
      {\textbf{\href{https://github.com/NoobMaster-version/automatic-animal-feeder}{Automatic Animal Feeder}} $|$ \emph{Servos, 3D Printing, Timer Control}}{2024}
      \resumeItemListStart
        \resumeItem{Built a scheduled dispenser that feeds the robotics lab cat through a 3D-printed chute when the lab is empty; still running, with an open-source release planned for animal shelters}
      \resumeItemListEnd

    \resumeProjectHeading
      {\textbf{River Cleaning Boat} $|$ \emph{RC Control, Mechanical Design, Waterproofing}}{2024}
      \resumeItemListStart
        \resumeItem{Built a remote-controlled boat with a front scoop and internal basket to collect floating trash, for an environmental studies course; the engineering problem was waterproofing electronics on a student budget}
      \resumeItemListEnd

  \resumeSubHeadingListEnd

%-----------COMPETITIONS-----------
\section{Competitions}
  \resumeSubHeadingListStart
    \resumeProjectHeading
      {\textbf{\href{https://noobmaster-version.github.io/spongefal/projects/eyantra.html}{eYantra Warehouse Drone Competition}} $|$ \emph{IIT Bombay}}{Oct 2024 -- Feb 2025}
      \resumeItemListStart
        \resumeItem{Developed autonomous navigation and PID control algorithms for stable hovering and motion of a simulated drone}
        \resumeItem{Implemented ArUco marker detection using computer vision for autonomous localization and task execution}
        \resumeItem{Reached the national semifinals; the fragility of camera-only localization here led directly to the UWB work at Krishna Defence}
      \resumeItemListEnd

    \resumeProjectHeading
      {\textbf{Mass Robotics Challenge} $|$ \emph{Quadruped Locomotion}}{2025}
      \resumeItemListStart
        \resumeItem{Entered Kurma, the turtle-inspired quadruped, advancing through to later competition rounds}
      \resumeItemListEnd
  \resumeSubHeadingListEnd

%-----------POSITIONS OF RESPONSIBILITY-----------
\section{Positions of Responsibility}
  \resumeSubHeadingListStart
    \resumeSubheading
      {Society of Robotics and Automation (SRA)}{Aug 2024 -- Mar 2026}
      {Mechanical Head \& Active Member}{VJTI, Mumbai}
      \resumeItemListStart
        \resumeItem{Conducted technical workshops on Embedded C, ESP32 microcontrollers, and ROS-based systems for 200+ students}
        \resumeItem{Mentored junior teams in developing self-balancing and line-following robots, providing hardware/software expertise}
        \resumeItem{Taught sensor fusion in practice, covering MPU6050 integration and complementary filtering for orientation estimation}
      \resumeItemListEnd

    \resumeSubheading
      {Pratibimb, VJTI Cultural Festival}{Aug 2023 -- Mar 2025}
      {Department Head, Electrical Engineering}{VJTI, Mumbai}
      \resumeItemListStart
        \resumeItem{Led the Electrical Engineering branch contingent across events, coordinating teams and participation}
      \resumeItemListEnd
  \resumeSubHeadingListEnd

%-----------ACHIEVEMENTS-----------
\section{Achievements}
 \begin{itemize}[leftmargin=0.3in, label=\textendash]
    \resumeItem{Led the Electrical Engineering branch to its first overall win at Pratibimb in the competition's 26-year history}
    \resumeItem{Directed \emph{Plot Se Panga}, a short film that won first prize across all branches at VJTI}
    \resumeItem{Reached the national semifinals of the eYantra Warehouse Drone Competition, IIT Bombay}
    \resumeItem{Placed 4th in the SRA maze-solving robot competition}
    \resumeItem{Consistent team performances across hackathons through first and second year at VJTI}
 \end{itemize}\vspace{-7pt}

%-----------TECHNICAL SKILLS-----------
\section{Technical Skills}
 \begin{itemize}[leftmargin=0.15in, label={}]
    \small{\item{
     \textbf{Programming Languages}{: C, C++, Python, Embedded C} \\
     \textbf{Software \& Frameworks}{: ROS2, Gazebo, RViz, ESP-IDF, Zephyr RTOS, STM32CubeIDE, KiCAD, micro-ROS} \\
     \textbf{Hardware}{: ESP32, STM32, Raspberry Pi, UWB Modules, Motor Drivers, BLDC/Stepper/DC Motors, LiDAR} \\
     \textbf{Core Competencies}{: Robot Kinematics, Field-Oriented Control, SLAM, NAV2, PID Control, PCB Design, IoT}
    }}
 \end{itemize}

\end{document}
